\documentclass[11pt]{article}

\usepackage[margin=1in]{geometry}
\usepackage{amsmath,amssymb,amsthm,mathtools}
\usepackage{hyperref}

\newcommand{\R}{\mathbb{R}}
\newcommand{\Sym}{\operatorname{Sym}}

\newtheorem{theorem}{Theorem}
\newtheorem{proposition}[theorem]{Proposition}
\newtheorem{corollary}[theorem]{Corollary}
\theoremstyle{remark}
\newtheorem{remark}[theorem]{Remark}
\newtheorem{example}[theorem]{Example}

\title{Boundary Persistence vs Functional Persistence}
\author{Research note generated in sandbox}
\date{2026-04-05}

\begin{document}
\maketitle

\section{Question}

Suppose the member-specific boundary fast state $q_{\alpha,t}\in\R^p$ satisfies
\begin{equation}
\dot q_{\alpha,t}=\mathcal C(\rho_t,q_{\alpha,t}),
\qquad
\alpha=1,\dots,K,
\label{eq:q-dyn}
\end{equation}
and the predictor is
\[
f_\alpha(x,t)=H(x;\rho_t,q_{\alpha,t}).
\]
If two members start from different boundary fast weights, do those differences necessarily disappear?

\medskip

There are in fact two distinct questions:
\begin{itemize}
\item Do the \emph{parameter differences} $q_{\alpha,t}-q_{\beta,t}$ disappear?
\item Do the resulting \emph{predictive differences} $f_\alpha(\cdot,t)-f_\beta(\cdot,t)$ disappear?
\end{itemize}

The answers need not coincide.

\section{Local Terminal Test}

Assume the collapsed mean-field trajectory converges to a stationary law $(\rho_\star,q_\star)$ with
\[
\mathcal C(\rho_\star,q_\star)=0.
\]
Write
\[
A_\star := \partial_q \mathcal C(\rho_\star,q_\star),
\qquad
J_\star(x) := \partial_q H(x;\rho_\star,q_\star).
\]
For two nearby members, define
\[
\Delta_t := q_{\alpha,t}-q_{\beta,t}.
\]
Then near the terminal law,
\begin{equation}
\dot \Delta_t = A_\star \Delta_t + o(\|\Delta_t\|),
\label{eq:delta-lin}
\end{equation}
and the predictive difference is
\begin{equation}
f_\alpha(x,t)-f_\beta(x,t)
=
J_\star(x)\Delta_t + o(\|\Delta_t\|).
\label{eq:f-lin}
\end{equation}

\begin{remark}
Equation \eqref{eq:delta-lin} governs whether the boundary parameters remain separated. Equation \eqref{eq:f-lin} governs whether that separation is visible at the level practitioners actually care about, namely prediction.
\end{remark}

\section{Parameter Collapse Is Not Automatic}

\begin{proposition}[Contractive terminal dynamics imply parameter collapse]
\label{prop:parameter-collapse}
If
\[
\Sym(A_\star)\preceq -\gamma I
\qquad
\text{for some }\gamma>0,
\]
then every sufficiently small terminal boundary difference decays exponentially:
\[
\|\Delta_t\|
\le
e^{-\gamma (t-t_0)}\|\Delta_{t_0}\|
\]
for all large enough $t\ge t_0$.
\end{proposition}

\begin{proof}
Ignoring the higher-order remainder in \eqref{eq:delta-lin}, the linearized energy satisfies
\[
\frac{d}{dt}\frac12\|\Delta_t\|^2
=
\langle \Delta_t,A_\star\Delta_t\rangle
\le
-\gamma \|\Delta_t\|^2.
\]
The nonlinear remainder is absorbed for sufficiently small $\|\Delta_t\|$.
\end{proof}

\begin{corollary}
Boundary fast weights do \emph{not} disappear by default. They disappear only in a dynamically contractive regime.
\end{corollary}

\section{Persistence of Parameters Need Not Mean Persistence of Prediction}

Let
\[
E_{\mathrm{ns}}
:=
\bigoplus_{\Re \lambda \ge 0} E_\lambda(A_\star)
\]
denote the non-stable generalized eigenspace of $A_\star$.

\begin{proposition}[Predictively silent persistence]
\label{prop:silent-persistence}
Assume there exists a nonzero subspace $E\subseteq E_{\mathrm{ns}}$ such that
\[
J_\star(x) v = 0
\qquad
\text{for every }v\in E \text{ and every }x.
\]
Then boundary differences initialized in $E$ may persist to first order in parameter space without producing any first-order predictive disagreement.
\end{proposition}

\begin{proof}
If $\Delta_{t_0}\in E$, then the linearized flow $e^{A_\star (t-t_0)}\Delta_{t_0}$ remains in $E$ because $E$ is $A_\star$-invariant. Since $E\subseteq E_{\mathrm{ns}}$, the linear dynamics does not force exponential decay along $E$. On the other hand, \eqref{eq:f-lin} vanishes identically on $E$ because $J_\star(x)E=\{0\}$ for every $x$.
\end{proof}

\begin{remark}
So ``boundary fast weights do not disappear'' is still not enough for useful diversity. The surviving directions must also be visible through the predictor Jacobian.
\end{remark}

\section{Visible Persistence Gives Real Terminal Diversity}

\begin{proposition}[Visible non-stable directions create terminal disagreement]
\label{prop:visible-persistence}
Assume there exists a vector $v\neq 0$ such that
\[
v\in E_{\mathrm{ns}}
\qquad\text{and}\qquad
\int \|J_\star(x)v\|^2\,\mu(\mathrm dx) > 0
\]
for the evaluation distribution $\mu$.
Then arbitrarily small initial boundary differences with a component along $v$ produce non-vanishing first-order predictive disagreement near the terminal law.
\end{proposition}

\begin{proof}
Choose $\Delta_{t_0}=\varepsilon v$. Since $v\in E_{\mathrm{ns}}$, the linearized terminal flow does not drive the coefficient of $v$ to zero exponentially. Therefore the first-order term in \eqref{eq:f-lin} retains a nonzero component proportional to $J_\star(x)v$. The stated integral condition guarantees that this component is nontrivial in prediction space.
\end{proof}

\section{A Better Practical Criterion}

The correct practical question is therefore not
\begin{quote}
Do boundary fast weights disappear?
\end{quote}
but rather
\begin{quote}
Which boundary directions are both dynamically persistent and predictively visible?
\end{quote}

This gives a two-part test at the terminal law:
\begin{itemize}
\item \textbf{Dynamics test.}
Check the terminal Jacobian $A_\star$. Directions with strongly negative decay rates are transient.
\item \textbf{Visibility test.}
Check the stationary predictor Gramian
\[
W_\star
:=
\int J_\star(x)^\top J_\star(x)\,\mu(\mathrm dx).
\]
Directions in or near $\ker W_\star$ are predictively silent.
\end{itemize}

\begin{corollary}[Safe design principle]
\label{cor:safe-principle}
Boundary fast weights are a meaningful place to encode finite-$K$ diversity only in directions that are simultaneously
\begin{itemize}
\item weakly decaying or persistent under $A_\star$, and
\item non-silent under $W_\star$.
\end{itemize}
\end{corollary}

\section{Examples}

\begin{example}[Persistent but useless]
Suppose $\mathcal C(\rho,q)=b(\rho)$ is independent of $q$. Then boundary differences persist exactly. But if the predictor also satisfies $H(x;\rho,q)=H(x;\rho)$ near the terminal law, then $J_\star(x)=0$ and the persistence is functionally irrelevant.
\end{example}

\begin{example}[Persistent and useful]
If $A_\star$ has a neutral direction $v$ and $J_\star(x)v\not\equiv 0$, then a small boundary perturbation along $v$ yields a non-vanishing terminal disagreement to first order.
\end{example}

\section{Takeaway}

The user's intuition is correct, but it needs one refinement:
\begin{quote}
Different boundary fast weights do not necessarily disappear. However, what matters for finite-$K$ uncertainty is not mere parameter persistence, but \emph{predictively visible persistence}.
\end{quote}

So the right end-of-training objects are the pair $(A_\star,W_\star)$, not just the existence of different initial boundary values.

\end{document}
